\documentclass[journal]{IEEEtran}

\usepackage{graphicx}
\usepackage{listings}
\usepackage{url}
\usepackage{booktabs}
\usepackage{multirow}
\usepackage{amsmath}

\title{A Modular Cyber-Physical Risk Assessment Architecture Integrating SIEM Telemetry and P\&ID-Derived System Graphs}

\author{Author Name\\
Affiliation\\
Email}

\begin{document}
\maketitle

\begin{abstract}
This paper presents a modular architecture for cyber-physical risk assessment that unifies
SIEM telemetry, host-based evidence, and process-level system graphs derived from P\&ID
artifacts. The system integrates a data-collection agent, a SIEM pipeline, a risk assessment
service, and a diagram-to-graph knowledge extraction workflow. We detail each component,
its interfaces, data schemas, and end-to-end orchestration, enabling reproducible analysis
across cyber and physical domains.
\end{abstract}

\begin{IEEEkeywords}
SIEM, cyber-physical systems, risk assessment, P\&ID, system graph, telemetry, OpenSearch
\end{IEEEkeywords}

\section{Introduction}
Cyber-physical systems (CPS) blend networked digital control with physical processes.
Assessing risk in such environments requires synchronized views of cyber telemetry and
process dependencies. This work describes an architecture that operationalizes this view by
integrating a SIEM pipeline with a process graph derived from P\&ID diagrams and system
models. The result is a modular framework that supports experimentation, telemetry ingestion,
system modeling, and risk analysis.

\section{System Overview}
Figure~\ref{fig:overview} summarizes the architecture. The system is composed of:
(1) a host agent for telemetry and evidence collection, (2) a SIEM pipeline for normalization
and search, (3) a P\&ID knowledge extraction pipeline that generates a process graph, and
(4) an ICS risk assessment service that consumes the graph and evidence to compute risk.

\begin{figure}[t]
  \centering
  \includegraphics[width=0.95\columnwidth]{figures/overview.pdf}
  \caption{High-level architecture. Telemetry flows from host agents into the SIEM pipeline and
  OpenSearch. P\&ID artifacts are converted to a system graph used by the risk assessment service.
  The dashboard orchestrates ingestion, visualization, and analysis.}
  \label{fig:overview}
\end{figure}

\section{Data Model and Schemas}
\subsection{System Graph (sim\_system.json)}
The process graph is represented as a JSON document with two top-level keys:
\texttt{variables} and \texttt{connections}. Variables are keyed by ID and include metadata
(e.g., type, module, domain). Connections are edges describing flow between variables.

\subsection{SIEM Event Schema}
Events are normalized to a consistent schema for indexing and querying. This includes
source metadata, timestamps, event type, and optional enrichment fields. The schema
supports extension for research instrumentation and correlation.

\section{P\&ID Knowledge Extraction}
P\&ID artifacts authored in draw.io are exported as XML and parsed into the system graph.
The parser supports both compressed and uncompressed draw.io XML. Nodes are derived
from shapes with metadata stored in draw.io ``Edit Data'' fields. Edges are derived from
connectors and converted into flow connections. Cyber vs. physical domain classification
is supported via explicit \texttt{domain} fields or keyword inference.

\subsection{Workflow}
\begin{enumerate}
  \item Author or import a draw.io P\&ID diagram.
  \item Export uncompressed XML.
  \item Convert XML to \texttt{sim\_system.json}.
  \item Optionally copy the generated file into the risk assessment workspace.
\end{enumerate}

\section{Host Agent and Telemetry Collection}
The host agent collects system telemetry, SBOMs, file integrity monitoring data, and
osquery results. It supports a command channel for ad-hoc data collection and posts
structured artifacts to server endpoints. These artifacts are normalized and ingested into
the SIEM pipeline for indexing.

\section{SIEM Pipeline}
The SIEM pipeline ingests events from host agents and external feeds, normalizes records
into a unified schema, and supports search and correlation. Events are stored in OpenSearch
for scalable query performance and dashboards.

\subsection{OpenSearch Integration}
OpenSearch is deployed as a single-node cluster for local development. Events are indexed
with a time-partitioned index naming scheme and managed by an index lifecycle policy.

\section{Risk Assessment Service}
The risk assessment service consumes the system graph and evidence to produce
probabilistic risk outputs. It exposes endpoints for status, node enumeration, and
probability computation. This service can be deployed alongside the SIEM pipeline and
accessed through the dashboard.

\subsection{Risk Service Workflow}
The risk service operates over a Dynamic Bayesian Network (DBN), which models time-indexed
dependencies between cyber and physical variables. The DBN is generated offline from the
system graph and stored as a BIFXML artifact. The runtime API exposes a thin layer over
this model for probability queries and vulnerability-driven updates.

The end-to-end workflow is: (1) upload a consolidated system graph
(\texttt{sim\_system.json}) to the service; (2) verify DBN availability via the status
endpoint; (3) issue evidence queries for probability distributions; and (4) submit cyber
vulnerability data to update the DBN and recompute probabilities.

\subsection{Integration Smoke Test}
To validate the integration, a lightweight smoke test is executed from the dashboard
codebase. The test posts \texttt{sim\_system.json}, queries \texttt{/status} and
\texttt{/nodes}, runs a minimal \texttt{/evidence} query on a small node sample, and
optionally exercises the \texttt{/probability} endpoint. This provides a reproducible
sanity check for the end-to-end data path before running full-scale experiments.

\section{Dashboard Orchestration and Visualization}
The dashboard coordinates ingestion and analysis. It provides:
(1) a P\&ID import UI for conversion and upload, (2) a system view with cyber/physical
visual differentiation and network overlays, and (3) risk assessment workflows to map
nodes, compute probabilities, and visualize results.

\section{Security Considerations}
The architecture includes token-based ingestion controls, optional TLS verification,
and role-based SIEM access. Default deployments disable OpenSearch security for local
experimentation but can be hardened for production environments.

\section{Evaluation Plan}
We plan to evaluate the architecture on (1) ingestion throughput, (2) graph generation
fidelity, (3) end-to-end latency from telemetry to risk outputs, and (4) analyst usability.

\section{Limitations and Future Work}
Current limitations include coarse mapping between cyber telemetry and physical variables,
and the absence of automated DBN regeneration on graph updates. Future work includes
model calibration, automated validation against ground truth, and multi-sensor fusion.

\section{Reproducibility}
All pipelines are deployed via Docker Compose, and configuration is controlled by
environment variables. The knowledge extraction pipeline is deterministic given a
P\&ID XML input.

\section{Conclusion}
This architecture provides a practical foundation for CPS risk analysis by unifying
SIEM telemetry with process-level system graphs. It supports reproducible research and
can be extended for broader instrumentation and analytical tasks.

\section*{Acknowledgments}
Optional.

\bibliographystyle{IEEEtran}
\bibliography{references}

\end{document}
